\documentclass[11pt]{article}
\usepackage[margin=1in]{geometry}
\usepackage{amsmath,amssymb,amsthm,mathtools,graphicx,booktabs}
\usepackage[hidelinks]{hyperref}
\newtheorem{theorem}{Theorem}
\newtheorem{remark}{Remark}
\newcommand{\R}{\mathbb R}
\newcommand{\qtilde}{\widetilde q}
\newcommand{\CMQd}{\mathrm{CMQd}}
\setlength{\parindent}{0pt}
\setlength{\parskip}{0.45em}

\title{A One-Dimensional Counterexample to\\
\texorpdfstring{$G_M=F_M$}{G M = F M} for CMQd Sets}
\author{Run \texttt{fa\_banach\_001}, lane 10}
\date{11 August 2026}

\begin{document}
\maketitle

\textbf{Status.} Candidate full counterexample, likely valid, for expert
review.

\section{Source question}

Let $E$ be a real Banach space, $B=E\times E^*$, and
\[
 q(x,x^*)=\langle x,x^*\rangle,\qquad
 r=q+\tfrac12\|\cdot\|^2 .
\]
The source calls a subset $M\subset B$ \emph{CMQd} when it is closed,
monotone, and quasidense, the last condition meaning
\[
 \inf_{m\in M}r(m-b)=0\qquad(b\in B).
\]
Writing $L(x,x^*)=(x^*,\widehat x)$ and
$\qtilde(y^*,y^{**})=\langle y^*,y^{**}\rangle$, it defines, for a
maximally monotone $M$,
\begin{align*}
 P_M(b)&=-\inf_{m\in M}q(m-b),\\
 G_M(b^*)&=-\inf_{m\in M}\qtilde(Lm-b^*),\\
 F_M(b^*)&=-\inf_{b\in B}
   \bigl[P_M(b)+\qtilde(b^*-Lb)\bigr].
\end{align*}
Every CMQd set is maximally monotone
\cite[Theorem 4.2]{Simons}.  The source proves $G_M\leq F_M$ and asks in
Remark~11.4 whether equality holds for every CMQd set
\cite[Remark 11.4, PDF p.~20]{Simons}.

\begin{center}
 \includegraphics[width=\textwidth]{figures/open_problem_crop.png}
\end{center}

\section{Claimed contribution}

\begin{theorem}\label{thm:main}
Let $E=\R$ with its usual norm.  For every $\lambda>0$, the line
\[
 M_\lambda=\{(s,\lambda s):s\in\R\}\subset B=\R^2
\]
is CMQd.  Under the canonical identifications $B^*=\R^2$ and
$L(u,v)=(v,u)$, its three source-paper functions are
\begin{align}
 P_{M_\lambda}(u,v)
   &=\frac{(v-\lambda u)^2}{4\lambda},\label{eq:P}\\
 G_{M_\lambda}(y,z)
   &=\frac{(y-\lambda z)^2}{4\lambda},\label{eq:G}\\
 F_{M_\lambda}(y,z)
   &=
   \begin{cases}
    0,&y=\lambda z,\\
    +\infty,&y\neq\lambda z.
   \end{cases}\label{eq:F}
\end{align}
Consequently $G_{M_\lambda}\neq F_{M_\lambda}$.  In particular, for
$\lambda=1$ and $b^*=(0,1)$,
\[
 G_{M_1}(0,1)=\frac14
 \qquad\text{but}\qquad
 F_{M_1}(0,1)=+\infty.
\]
Thus the conjecture in Remark~11.4 is false, already in dimension one.
\end{theorem}

\section{Proof intuition}

The penalty $P_{M_\lambda}(u,v)$ controls only the transverse coordinate
$v-\lambda u$.  In the definition of $F_M$, the remaining quadratic
expression instead contains $(\lambda u+v)^2$.  It therefore has a null
direction $(u,v)=(t,-\lambda t)$.  Unless the dual point lies on the line
$y=\lambda z$, the linear part drifts to $-\infty$ along this direction,
forcing $F_M=+\infty$.  Direct minimization over the graph defining $G_M$
has no such escape and always gives a finite quadratic value.

\section{Proof of Theorem~\ref{thm:main}}

Fix $\lambda>0$.  The set $M_\lambda$ is closed.  It is monotone because
\[
 q\bigl((s,\lambda s)-(t,\lambda t)\bigr)
   =\lambda(s-t)^2\geq0.
\]
For an arbitrary $b=(a,c)\in B$,
\begin{align*}
 r\bigl((s,\lambda s)-(a,c)\bigr)
 &= (s-a)(\lambda s-c)
    +\frac12\bigl[(s-a)^2+(\lambda s-c)^2\bigr]\\
 &=\frac12\bigl[(1+\lambda)s-(a+c)\bigr]^2.
\end{align*}
Its infimum over $s\in\R$ is zero, attained at
$s=(a+c)/(1+\lambda)$.  Hence $M_\lambda$ is quasidense and therefore
CMQd.

For $(u,v)\in B$, completing the square gives
\begin{align*}
 (s-u)(\lambda s-v)
 &=
 \lambda\left(s-\frac{v+\lambda u}{2\lambda}\right)^2
 -\frac{(v-\lambda u)^2}{4\lambda}.
\end{align*}
Taking the negative infimum proves \eqref{eq:P}.

Now let $b^*=(y,z)\in B^*$.  Since
$L(s,\lambda s)=(\lambda s,s)$,
\begin{align*}
 (\lambda s-y)(s-z)
 &=
 \lambda\left(s-\frac{y+\lambda z}{2\lambda}\right)^2
 -\frac{(y-\lambda z)^2}{4\lambda}.
\end{align*}
The definition of $G_M$ immediately gives \eqref{eq:G}.

It remains to compute $F_M$.  For $b=(u,v)$, the expression whose
infimum occurs in its definition is
\begin{align}
 J(u,v)
 &:=
 P_{M_\lambda}(u,v)+\qtilde\bigl((y,z)-L(u,v)\bigr)\notag\\
 &=\frac{(v-\lambda u)^2}{4\lambda}
   +(y-v)(z-u)\notag\\
 &=\frac{(\lambda u+v)^2}{4\lambda}
   -yu-zv+yz.\label{eq:J}
\end{align}
Along $(u,v)=(t,-\lambda t)$ this becomes
\[
 J(t,-\lambda t)=(\lambda z-y)t+yz.
\]
If $y\neq\lambda z$, this affine function is unbounded below, so
$\inf J=-\infty$ and $F_{M_\lambda}(y,z)=+\infty$.

If $y=\lambda z$, put $w=\lambda u+v$ in \eqref{eq:J}.  Then
\[
 J(u,v)
 =\frac{w^2}{4\lambda}-zw+\lambda z^2
 =\frac{(w-2\lambda z)^2}{4\lambda}.
\]
Its infimum is zero.  Thus $F_{M_\lambda}(y,z)=0$, proving
\eqref{eq:F} and the theorem.

\section{Stress checks and scope}

The value $+\infty$ in \eqref{eq:F} is permitted by the source's stated
codomain $F_M\colon B^*\to]-\infty,\infty]$.  No nonreflexive phenomenon,
closure subtlety, or unattained infimum is involved.  For the displayed
witness, one can check the divergence without any parameter algebra:
\[
 P_{M_1}(u,v)+(0-v)(1-u)=\frac{(u+v)^2}{4}-v.
\]
At $(u,v)=(-N,N)$ this equals $-N$, whereas the direct computation of
$G_{M_1}(0,1)$ is $-\inf_s s(s-1)=1/4$.

The theorem gives a complete description of $P_M,G_M,F_M$ for every
positive scalar graph, rather than merely one isolated witness.  It does
not attempt to characterize all CMQd subsets on which equality holds.

The script \texttt{code/verify\_counterexample.py} independently checks
every completed-square identity and the explicit witness in exact symbolic
arithmetic.  The analytic proof above does not depend on computation.

A bounded novelty check was performed on 11 August 2026.  The run's cheap
indexes were searched by arXiv id, exact conjecture text, title, CMQd,
quasidensity, and the pair $G_M,F_M$; the local citing corpus and targeted
web/arXiv searches were also checked.  No later proof, counterexample, or
matching scalar computation was located.  This is bounded negative
evidence, not a priority claim.

\textbf{Human-review recommendation.}  The result is elementary enough for
direct line-by-line review.  The critical checks are that the source uses
the displayed definition of $F_M$ with its leading minus sign, and that
the codomain explicitly allows $+\infty$.

\begin{thebibliography}{9}
\small
\bibitem{Simons}
Stephen Simons,
\emph{Faces of quasidensity},
arXiv:2407.00479 (2024),
\url{https://arxiv.org/abs/2407.00479}.
\end{thebibliography}

\end{document}
